\section{The objective, its logical meaning, and the comparison}

\subsection{What is being constructed}
Write $\NN=\{0,1,2,\ldots\}$. Fix an effective standard presentation of ZFC,
and let $\Con(\ZFC)$ mean that there is no finite derivation of falsity in
that presentation. We construct fixed integer-coefficient polynomials
$P_{12}$ and $F_9$, with the following meanings:
\begin{align}
 \ZFC&\vdash
 \bigl(\nexists\boldsymbol{x}\in\PP^{12}:P_{12}(\boldsymbol{x})=0\bigr)
 \longleftrightarrow\Con(\ZFC),\label{eq:goal12}\\
 \ZFC&\vdash
 \bigl(\nexists\boldsymbol{v}\in\NN^9:F_9(\boldsymbol{v})=0\bigr)
 \longleftrightarrow\Con(\ZFC).\label{eq:goal9}
\end{align}
The coefficients are independent of whether ZFC is consistent. None is
chosen by waiting for a refutation or a halting computation. Every displayed
auxiliary letter in the polynomial definitions denotes a substituted
polynomial expression; it is not an extra quantified unknown.

The domains in \eqref{eq:goal12} and \eqref{eq:goal9} are part of the claims.
In particular, the twelve-variable formula must not simply be interpreted
over unrestricted integers. The nine-variable formula can be converted to
nine positive unknowns by substituting $v_j=x_j-1$; this preserves exact
degree but adds nine subtraction gates to the supplied circuit.

\subsection{A precise meaning of simpler}
For an explicitly represented polynomial or system, at least five measures
matter: the number and domain of unknowns, total degree, the number of
equations, the number of arithmetic operations in a shared straight-line
circuit, and the description or bit size of its coefficients. These measures
compete. Replacing a system by a sum of squares removes equations without
removing unknowns and generally doubles the maximum degree. Replacing
several square roots by a primitive element removes unknowns but increases
degree. Fixed integers are free inputs in the circuit model, although
computing or printing those integers may be expensive.

Our operation model counts every binary $+$, $-$ and $\times$, including
multiplication by a literal integer and squaring. Reuse of a previously
computed value is free. Comparisons in a system are free. All fixed powers
in the nine-variable expression are expanded into multiplication gates.
The cost of constructing a fixed coefficient is recorded separately.

There is no presently justified single global optimum to report. The
concrete achievements are two improved points in this collection of size
tradeoffs. In particular, neither result reduces the existing nine-variable
count, and the earlier 67-operation equation and low-degree lifted systems
retain advantages on their own measures.

\subsection{Comparison with the current repository}
The following data summarize the current positive twelve-variable products;
the baseline entries are from the pinned catalogue and the individual
packages \cite{catalogue,p17,p18,p19,p20,p21,p22,p23}.
\begin{center}
\begin{tabular}{lrrl}
\toprule
Construction & Degree & Operations & Main method\\
\midrule
17-primitive12 & 2601 & 167 & primitive element\\
18-jones9, direct form & 1609 & 151 & radical combination\\
19-wang943, direct form & 1257 & 139 & saturated radical reduction\\
20-root12 & 1145 & 123 & nested root separation\\
21-norm12 & 1353 & 116 & efficient norm circuit\\
22-monic12 & 1529 & 117 & monic polynomial merger\\
23-separating12 & 1369 & 144 & separating coefficients\\
\textbf{This report} & \textbf{1081} & \textbf{116} & smaller uniform weights\\
\bottomrule
\end{tabular}
\end{center}
Package 16 contributes a thirteen-positive-variable form of degree $515$ at
$84$ operations, and a separate unrestricted-integer reduction; it is not a
twelve-variable competitor. The new product improves the previous minimum
twelve-variable degree by $64$, or approximately $5.59\%$, and reduces the
degree of the former 116-operation product by $272$, or approximately
$20.10\%$. The new three-equation system retains the best stated count of
$112$ operations while reducing its maximum degree to $1072$.

For nine natural unknowns, put $\mathcal L=5^{33}$:
\begin{center}
\begin{tabular}{lrr}
\toprule
Construction & Exact degree & Operations\\
\midrule
18, improved majorants & $161844\mathcal L+484$ & $20\,531$\\
This report & $93264\mathcal L+288$ & $20\,547$\\
\bottomrule
\end{tabular}
\end{center}
The variable count is unchanged. The degree falls by approximately
$42.374138\%$, while the circuit grows by sixteen operations, about
$0.078\%$. The exact new degree is
\[
10\,857\,358\,574\,867\,248\,535\,156\,250\,288.
\]
The complete base coefficient tuple is taken from package 19 rather than
splicing the coefficients of different compilers. This reuse is a synthesis
of existing work, not a newly discovered coefficient-size record.

\subsection{The proof chain and what is inherited}
The literal machine is K215, with initial state $0$, an all-zero two-sided
tape, and halt state $215$. The repository proves its ordinary calibration
\cite{k215}:
\[
 \operatorname{Halt}(K215,\mathrm{blank})
 \longleftrightarrow
 \text{there exists a finite ZFC refutation}.
\]
The current Lean endpoint named
\texttt{KEndpoint215.blank\_halts\_iff\_zfc\_refutation} and the corresponding
Rocq endpoint provide separate formal evidence for the machine theorem.
We inspected the current Lean endpoint but did not rebuild the full
Lean/Rocq developments. The mathematical ZFC-provability claim uses the
ordinary proof, not an inference from the soundness of Lean.

Package 19 supplies the marked periodic Wang-tile language and the entire
fixed coefficient tuple. Its interface theorem identifies the resulting
native mask constraints with a blank halt. The inherited Pell arithmetic
recovers those constraints from its equations; conversely it constructs
canonical positive witnesses from a halting tableau. The signed main-root
and signed auxiliary-quotient extensions are required when roots are
eliminated. The proofs retained in \texttt{vendor/inherited/} record those
extensions. The new twelve-variable argument changes only the root-combining
part of that arithmetic proof.

The nine-variable construction starts from the concrete 32-positive-variable
quadratic system already specialized in package 18. Its sum of squares is
a quartic. Jones's compression and its repository implementation give an
equivalent condition with eight outer natural variables and four arithmetic
tests \cite{jones,p18,ninemain}. We replace the final polynomial combining
those tests and its root bounds. Both new steps receive complete proofs
below.

Thus the proof structure is a composition of stated theorems, with a clear
boundary between reused compiler mathematics, new algebra, and finite
computations. A checksum confirms which integer or machine table was used;
it does not prove the unbounded compiler theorem. A circuit identity confirms
an algebraic transformation; it does not establish consistency.

\subsection{Consequences and limitations of the equivalences}
The conclusions are representation theorems, not proofs of either consistency
or inconsistency. Under the usual consistency hypothesis for ZFC, G\"odel's
second incompleteness theorem prevents ZFC from proving the unsolvability
side of \eqref{eq:goal12} or \eqref{eq:goal9}. Consistency alone should not be
used to assert that ZFC cannot prove the opposite, existential statement;
an additional arithmetic soundness hypothesis is needed for that stronger
two-sided assertion. Our construction does not require either hypothesis
to prove the displayed biconditionals inside ZFC.
